\documentclass[11pt]{article}

\usepackage[margin=1in]{geometry}
\usepackage[T1]{fontenc}
\usepackage[utf8]{inputenc}
\usepackage{lmodern}
\usepackage{graphicx}
\usepackage{booktabs}
\usepackage{longtable}
\usepackage{array}
\usepackage{xcolor}
\usepackage[hidelinks]{hyperref}
\usepackage{enumitem}
\setlist{nosep}

\title{Guia de Estado Actual\\GhostMerc Frontier Territory}
\author{Cesar Josue Zapata Chilla}
\date{Marzo 2026}

\begin{document}

\maketitle
\tableofcontents
\newpage

\section{Que tienes ahora}

Ahora tienes un proyecto consolidado dentro de \texttt{PR\_LR\_H} con \textbf{Frontier Territory} como camino canonico. Ya no es una mezcla informal de demos: hay un entorno headless principal, entrenamiento PPO, evaluacion con reward oculto, renderer separado y una ruta inicial para GNN y preferencias.

\subsection{Pieza principal}
La pieza central del proyecto es:
\begin{itemize}
    \item \texttt{chromahack/envs/ghostmerc\_frontier\_env.py}
\end{itemize}

Ese entorno ya soporta dos modos de observacion:
\begin{itemize}
    \item \textbf{flat}: baseline MLP actual
    \item \textbf{dict}: ruta estructurada para el backend relacional/GNN
\end{itemize}

\subsection{Stack activo}
El stack que esta vivo y debes usar es:
\begin{itemize}
    \item \texttt{python -m chromahack.training.train\_ppo\_frontier}
    \item \texttt{python -m chromahack.evaluation.eval\_frontier\_hidden}
    \item \texttt{python -m chromahack.rendering.frontier\_dual\_renderer}
    \item \texttt{python -m chromahack.run\_experiment --mode frontier}
\end{itemize}

\subsection{Que ya no es el camino principal}
Quedan como legado:
\begin{itemize}
    \item bridge demo
    \item classic GhostMerc
    \item wrappers viejos de CNN / visual proxy
\end{itemize}

No estan borrados por completo, pero ya no forman parte de la historia principal ni del flujo recomendado.

\section{Arquitectura actual}

\subsection{Separacion de responsabilidades}
La arquitectura consolidada ya tiene esta forma:
\begin{enumerate}
    \item \textbf{Entorno Gymnasium headless}: contiene dinamica, rewards, detector de transicion y snapshots.
    \item \textbf{Entrenamiento PPO}: usa SB3 y entrena contra \texttt{proxy\_reward}.
    \item \textbf{Evaluacion oculta}: mide \texttt{true\_reward}, gap, exploit frequency y trayectorias.
    \item \textbf{Renderer PyGame}: solo consume trayectorias o una politica en inferencia.
    \item \textbf{Pipeline de preferencias}: exporta pares, entrena reward model y permite reentrenar con \texttt{pref\_model}.
    \item \textbf{Observabilidad}: registra entropia, deriva KL y una senal PCA simple.
\end{enumerate}

\subsection{Archivos importantes}
\begin{longtable}{p{0.37\textwidth}p{0.55\textwidth}}
\toprule
\textbf{Archivo} & \textbf{Rol} \\
\midrule
\endfirsthead
\toprule
\textbf{Archivo} & \textbf{Rol} \\
\midrule
\endhead
\texttt{chromahack/envs/ghostmerc\_frontier\_env.py} & entorno principal \\
\texttt{chromahack/envs/territory\_generator.py} & distritos, zonas y actores \\
\texttt{chromahack/training/train\_ppo\_frontier.py} & PPO canonico \\
\texttt{chromahack/evaluation/eval\_frontier\_hidden.py} & evaluacion con reward oculto \\
\texttt{chromahack/rendering/frontier\_dual\_renderer.py} & replay, showcase, master, trailer \\
\texttt{chromahack/gnn/feature\_extractor.py} & extractor relacional para \texttt{dict} obs \\
\texttt{chromahack/intervention/pref\_model.py} & export, train y wrapper de reward por preferencias \\
\texttt{chromahack/utils/metrics.py} & callback Frontier y agregados \\
\texttt{chromahack/metrics/frontier\_observability.py} & entropia, KL y PCA simple \\
\texttt{chromahack/run\_experiment.py} & orquestador baseline / rlhf \\
\bottomrule
\end{longtable}

\section{Que artefactos concretos tienes}

\subsection{Demo maestra larga}
La demo maestra canonica ahora esta dentro del repo:
\begin{itemize}
    \item \texttt{artifacts/demos/frontier\_master\_demo\_long/model/ppo\_best.zip}
    \item \texttt{artifacts/demos/frontier\_master\_demo\_long/eval\_frontier\_hidden/summary.json}
    \item \texttt{artifacts/demos/frontier\_master\_demo\_long/eval\_frontier\_hidden/trajectories/episode\_040.json}
\end{itemize}

\subsection{Resultados principales de esa demo}
En 50 episodios del distrito 5:
\begin{itemize}
    \item \texttt{avg\_proxy\_return = 816.61}
    \item \texttt{avg\_true\_return = -306.90}
    \item \texttt{proxy\_true\_gap = 1123.51}
    \item \texttt{avg\_armed\_neutral\_false\_positive\_rate = 0.713}
    \item \texttt{avg\_containment\_tick\_count = 84.96}
    \item \texttt{exploit\_frequency = 0.18}
    \item \texttt{phase\_transition\_rate = 0.18}
    \item \texttt{avg\_phase\_transition\_step = 100}
\end{itemize}

\subsection{Mejor episodio curado}
El episodio maestro actual es:
\begin{itemize}
    \item \texttt{episode\_040.json}
\end{itemize}

Sus senales mas llamativas:
\begin{itemize}
    \item \texttt{containment\_tick\_count = 485}
    \item \texttt{proxy\_true\_gap = 2039.68}
    \item \texttt{armed\_neutral\_false\_positive\_rate = 1.0}
    \item \texttt{phase\_transition\_step = 100}
    \item beat visual: \texttt{CONTAINMENT EXPLOIT ACTIVE}
\end{itemize}

\section{Que puedes ejecutar para ver ahora mismo}

\subsection{Importante}
El directorio recomendado es:
\begin{itemize}
    \item \texttt{C:\textbackslash Users\textbackslash pc\textbackslash Desktop\textbackslash Proyecto\_LR\textbackslash PR\_LR\_H}
\end{itemize}

Tambien deberia funcionar desde la carpeta padre gracias al shim, pero la raiz del repo es el lugar canonico.

\subsection{Para ver la demo maestra directamente}
\begin{verbatim}
python -m chromahack.rendering.frontier_dual_renderer master ^
  --demo_dir artifacts/demos/frontier_master_demo_long
\end{verbatim}

Eso abre la mejor trayectoria ya curada desde el resumen maestro.

\subsection{Para ver el episodio exacto}
\begin{verbatim}
python -m chromahack.rendering.frontier_dual_renderer replay ^
  --trajectory artifacts/demos/frontier_master_demo_long/
eval_frontier_hidden/trajectories/episode_040.json
\end{verbatim}

\subsection{Para verlo con presentacion mas clara}
\begin{verbatim}
python -m chromahack.rendering.frontier_dual_renderer showcase ^
  --trajectory artifacts/demos/frontier_master_demo_long/
eval_frontier_hidden/trajectories/episode_040.json ^
  --slow_mo 0.8 --pause_frames 24 --alert_linger_frames 120
\end{verbatim}

\subsection{Para ver el corte editorial}
\begin{verbatim}
python -m chromahack.rendering.frontier_dual_renderer trailer_cut ^
  --demo_dir artifacts/demos/frontier_master_demo_long
\end{verbatim}

\subsection{Para exportar video}
\begin{verbatim}
python -m chromahack.rendering.frontier_dual_renderer trailer_cut ^
  --demo_dir artifacts/demos/frontier_master_demo_long ^
  --export_dir artifacts/captures/frontier_trailer_cut_final ^
  --video_name ghostmerc_frontier_final.mp4
\end{verbatim}

\subsection{Para ver una politica viva}
\begin{verbatim}
python -m chromahack.rendering.frontier_dual_renderer policy ^
  --model_dir artifacts/demos/frontier_master_demo_long/model ^
  --model_name ppo_best ^
  --district_id 5 ^
  --camera_mode cinematic ^
  --alert_linger_frames 120 ^
  --auto_reset --max_episodes 12
\end{verbatim}

\section{Que puedes ejecutar para entrenar}

\subsection{Baseline MLP}
\begin{verbatim}
python -m chromahack.training.train_ppo_frontier ^
  --out_dir artifacts/models/frontier_baseline ^
  --total_steps 200000 ^
  --n_envs 4 ^
  --policy_backend mlp ^
  --observation_mode flat ^
  --reward_mode proxy
\end{verbatim}

\subsection{Evaluacion del baseline}
\begin{verbatim}
python -m chromahack.evaluation.eval_frontier_hidden ^
  --model_dir artifacts/models/frontier_baseline ^
  --model_name ppo_best ^
  --n_episodes 50 ^
  --save_trajectories
\end{verbatim}

\subsection{GNN / backend relacional}
\begin{verbatim}
python -m chromahack.training.train_ppo_frontier ^
  --out_dir artifacts/models/frontier_gnn ^
  --total_steps 200000 ^
  --n_envs 4 ^
  --policy_backend gnn ^
  --observation_mode dict
\end{verbatim}

\subsection{Master demo largo}
\begin{verbatim}
python -m chromahack.run_experiment --mode frontier --phase baseline --master_demo
\end{verbatim}

\section{Que puedes ejecutar para la ruta RLHF}

\subsection{Pipeline completo rapido}
\begin{verbatim}
python -m chromahack.run_experiment --mode frontier --phase rlhf --quick
\end{verbatim}

Ese comando ya fue validado en esta iteracion y hace:
\begin{enumerate}
    \item baseline corto,
    \item export de pares de preferencia,
    \item entrenamiento del reward model,
    \item reentrenamiento corto con \texttt{pref\_model},
    \item evaluacion del run reentrenado.
\end{enumerate}

\subsection{Pipeline paso a paso}
\begin{verbatim}
python -m chromahack.intervention.pref_model export ^
  --trajectory_dir artifacts/frontier/baseline/eval_frontier_hidden/trajectories ^
  --out_dir artifacts/frontier/preferences

python -m chromahack.intervention.pref_model train ^
  --dataset artifacts/frontier/preferences/preference_pairs.json ^
  --out_dir artifacts/frontier/preferences/model

python -m chromahack.training.train_ppo_frontier ^
  --out_dir artifacts/frontier/pref_model_run/model ^
  --reward_mode pref_model ^
  --reward_model_path artifacts/frontier/preferences/model/preference_reward_model.pt ^
  --policy_backend mlp ^
  --observation_mode flat
\end{verbatim}

\subsection{Estado actual de la ruta de preferencias}
La ruta ya esta funcional, no solo en scaffolding. En la ultima validacion:
\begin{itemize}
    \item \texttt{train\_accuracy = 0.80}
    \item \texttt{validation\_accuracy = 0.7647}
\end{itemize}

Eso confirma que el reward model sintetico ya se entrena y se puede inyectar en el entorno.

\section{Que quedo validado tecnicamente}

\subsection{CLI}
Ya funcionan:
\begin{itemize}
    \item \texttt{python -m chromahack.training.train\_ppo\_frontier --help}
    \item \texttt{python -m chromahack.evaluation.eval\_frontier\_hidden --help}
    \item \texttt{python -m chromahack.run\_experiment --help}
    \item \texttt{python -m chromahack.intervention.pref\_model --help}
    \item \texttt{python -m chromahack.rendering.frontier\_dual\_renderer --help}
    \item \texttt{python -m chromahack.rendering.frontier\_dual\_renderer master --help}
\end{itemize}

\subsection{Tests}
Quedo validado:
\begin{itemize}
    \item \texttt{python scripts/smoke\_test.py}
    \item \texttt{python -m pytest tests/test\_frontier\_observation\_modes.py}
\end{itemize}

Las pruebas locales cubren:
\begin{itemize}
    \item paridad de dinamica entre \texttt{flat} y \texttt{dict}
    \item wrapper de reward por preferencias
\end{itemize}

\subsection{Validaciones funcionales extra}
Tambien quedo probado:
\begin{itemize}
    \item entrenamiento corto con backend \texttt{gnn}
    \item evaluacion de un checkpoint \texttt{gnn}
    \item export y entrenamiento del reward model
    \item entrenamiento corto con \texttt{reward\_mode=pref\_model}
    \item orquestador completo \texttt{--phase rlhf --quick}
\end{itemize}

\section{Como interpretar lo que ves}

\subsection{Proxy reward}
Es lo que el agente optimiza durante entrenamiento baseline.

\subsection{True reward}
Es la evaluacion oculta que si representa el objetivo real.

\subsection{Proxy-true gap}
Si sube mucho, hay divergencia entre score visible y exito real.

\subsection{Armed neutral false positive rate}
Si sube, el agente esta clasificando mal a actores armados que no deberia tratar como amenaza real.

\subsection{Containment tick count}
Si sube, el agente esta cobrando puntos por mantener actores heridos y vigilados en vez de resolver la situacion.

\subsection{Phase transition}
Si aparece un \texttt{phase\_transition\_step}, ese es el paso donde la conducta ya entro a un regimen diferente y mas explotativo.

\section{Resumen rapido para explicarlo facil}

\subsection{En una frase}
Entrene un agente tactico con una recompensa proxy corrupta, y el proyecto muestra como puede subir mucho su score mientras fracasa en el objetivo real.

\subsection{En 30 segundos}
GhostMerc Frontier Territory es un testbed de RL para estudiar reward hacking. El agente vive en un distrito con civiles, hostiles y actores ambiguos. PPO lo entrena con una senal proxy que premia headshots, threat tags y containment ticks. El objetivo real oculto es proteger civiles, estabilizar el territorio y evitar falsos positivos. En la demo larga, el agente logra proxy muy alto y true reward muy negativo, lo que demuestra desalineacion visible y medible.

\subsection{Tres numeros que conviene memorizar}
\begin{itemize}
    \item \texttt{816.61}: proxy medio de la demo larga
    \item \texttt{-306.90}: true reward medio de la demo larga
    \item \texttt{2039.68}: gap del mejor episodio
\end{itemize}

\section{Que hacer despues}

Si lo que quieres es \textbf{ver} el proyecto, ejecuta:
\begin{itemize}
    \item \texttt{master}
    \item \texttt{showcase}
    \item \texttt{trailer\_cut}
\end{itemize}

Si lo que quieres es \textbf{investigarlo}, ejecuta:
\begin{itemize}
    \item baseline MLP
    \item backend GNN
    \item fase RLHF sintetica
\end{itemize}

Si lo que quieres es \textbf{explicarlo facil}, muestra:
\begin{itemize}
    \item el replay del \texttt{episode\_040}
    \item el gap proxy--true
    \item la idea de que el agente parece competente pero optimiza el objetivo equivocado
\end{itemize}

\section{Cierre}

En este momento tu proyecto ya no es solo una idea o una demo visual. Ya tienes:
\begin{itemize}
    \item entorno Frontier consolidado
    \item baseline reproducible
    \item renderer utilizable
    \item demo maestra larga dentro del repo
    \item backend relacional habilitado
    \item ruta inicial RLHF funcionando
    \item documentacion y paper separados
\end{itemize}

Lo mas importante: ya puedes \textbf{mostrar}, \textbf{explicar} y \textbf{seguir investigando} sobre el mismo stack, sin rehacer el proyecto desde cero.

\end{document}
